% ===================== Section I: Introduction (Member 1) =====================
\section{Introduction}
\label{sec:intro}

Opening a bank account used to mean going to a branch with an identity card. Now most banks and
mobile wallets let people apply from their phone in a few minutes, and most everyday payments have
moved to digital channels. In Pakistan, the State Bank of Pakistan reports that 88\% of retail
transactions in FY2024--25 were digital, and that payments made through banking, branchless banking
and e-money wallet apps grew by 52\% to about 6.2 billion transactions in the same year
\cite{sbp2025payment}. Fraudsters use this convenience too. In account opening fraud, someone
applies with a stolen or made-up identity, gets the account approved, and then uses it to receive
illegal money or to spend the credit line that comes with it. Since the real owner of the identity
cannot be traced, the bank usually carries the loss \cite{jesus2022turning}. The problem is not a
small one: in 2024 the U.S. Federal Trade Commission received more than 1.1 million identity theft
reports, and consumers reported losing over \$12.5 billion to fraud \cite{ftc2025sentinel}.

To deal with this, banks score every new application with a machine learning (ML) model and send
the riskiest ones for rejection or manual review. Building such a model is harder than a normal
classification task, for three reasons. The first is that fraud is rare. In our data only about one
application in a hundred is fraud, so a model that never flags anyone is already 99\% accurate while
catching nothing. The second is that fraud changes over time. Fraudsters adapt to the checks they
run into, and genuine customers also change from month to month, so a model that does well on its
training months can be weaker on the months that follow \cite{dalpozzolo2018credit}. The third is
that the decision affects people. A false positive here means an honest customer is refused a bank
account, and if this happens more often to one age group than to another, the harm is not shared
equally \cite{jesus2022turning}.

A lot of the published work on fraud detection does not take these problems into account. Many
studies use public credit card data and report accuracy as their main result, even though the
classes are extremely imbalanced \cite{alarfaj2022credit, taha2020intelligent}. Others split the
data randomly and classify at the default 0.5 threshold, and a recent study on bank account opening
data did both \cite{sun2025objective}. A random split mixes later months into the training data,
which hides any drop caused by change over time. A fixed 0.5 threshold also does not match how a
bank works in practice, since it can only afford a limited number of false alarms. The Bank Account
Fraud (BAF) suite released by Jesus~\textit{et al.}~\cite{jesus2022turning} makes a more realistic
study possible. It has one million synthetic but realistic applications over eight months, a month
column for time-based evaluation, and protected attributes such as age. The original paper trained
LightGBM models on the first six months, tested them on the last two, and reported the true positive
rate at a 5\% false positive rate together with a fairness measure between age groups. It did,
however, set the 5\% threshold on the test months themselves, left a separate validation split for
the threshold as future work, and did not compare different model families under the same protocol.

In this paper we compare a majority-class baseline, a logistic regression scorecard and four ML
models from two families (random forest, LightGBM, CatBoost and a multilayer perceptron) on the BAF
Base dataset, using a strictly time-based protocol. We try to answer three research questions:
\begin{itemize}
  \item \textbf{RQ1:} Which model catches the most fraud at a fixed 5\% false positive rate?
  \item \textbf{RQ2:} How much does performance drop from the training months to the two test months?
  \item \textbf{RQ3:} Does the best model flag applicants over 50 at a different false positive rate
        from younger applicants, and what does it cost to correct this?
\end{itemize}

Our main contributions are:
\begin{enumerate}
  \item A time-based evaluation protocol for BAF without leakage. Hyperparameters are tuned on five
        forward-chaining folds over months 0--5, the 5\% FPR threshold is fixed on month 5, and
        months 6 and 7 are used only once, for the final test.
  \item A comparison of six models with the same folds, preprocessing, tuning budget and metrics,
        reported as mean $\pm$ standard deviation over the folds and month by month on the test
        period.
  \item A measurement of performance drift, including how far the false positive rate moves from
        its 5\% target when a threshold chosen on earlier data is used on later months.
  \item A fairness analysis of the false positive rate ratio between applicants over 50 and the
        rest, with the cost in recall of two corrections (group-specific thresholds and reweighing).
  \item An ablation study over six feature groups and two preprocessing steps, an error analysis of
        missed fraud cases, and a public repository that reproduces every number in the paper.
\end{enumerate}

The rest of the paper is organised as follows. Section~\ref{sec:related} reviews related work on
fraud detection, concept drift and algorithmic fairness. Section~\ref{sec:dataset} describes the
dataset, and Section~\ref{sec:method} explains the time-based validation and the rest of the
methodology. Section~\ref{sec:results} gives the results, Section~\ref{sec:fairness} the fairness
analysis and Section~\ref{sec:ablation} the ablation and error analysis.
Section~\ref{sec:discussion} discusses what we found and its limits, and
Section~\ref{sec:conclusion} concludes.
% NOTE for the group: check the \label names match the ones used in your sections.
